\documentclass[../../../include/open-logic-section]{subfiles}

\begin{document}

\olfileid{sth}{z}{pairs}
\olsection{Pasangan}

Aksioma sabanjure kang arep kita timbang yaiku:

\begin{axiom}[Pasangan]
Kanggo sembarang himpunan $a, b$, himpunan $\{a, b\}$ ana.
\[
	\forall a \forall b \exists P \forall x (x \in P \liff (x = a \lor x = b))
\]
\end{axiom}

Mangkene carane ndhukung aksioma iki nganggo konsepsi
iteratif. Upamane $a$ wis kasedhiya ing tataran $S$ lan $b$
wis kasedhiya ing tataran $T$. Upamane $M$ iku tataran kang
luwih pungkasan saka $S$ lan $T$. Amarga $a$ lan $b$
loro-lorone wis kasedhiya ing tataran $M$, himpunan $\{a,b\}$
minangka koleksi kang bisa kawangun ing sembarang tataran
sawise $M$.

Nanging enteni!{} Yagene kita nganggep \emph{ana} tataran
sawise $M$? Yen ora ana, panjelasan kita bakal gagal.
Mula kanggo ndhukung Pasangan, kita kudu nambah prinsip
liyane ing gambaran saka \olref[sth][z][story]{sec}, yaiku:
\begin{enumerate}
	\item[] \stagessucc. Ora ana tataran pungkasan.
\end{enumerate}
Apa prinsip iki duwe dhasar? Ora ana ing gambaran Shoenfield
kang nyebut kanthi \emph{gamblang} yen ora ana tataran
pungkasan. Nanging sanajan (sacara trep) iki tambahan
tumrap gambaran sadurunge, prinsip iki cocog karo gagasan
dhasar yen himpunan kawangun ing tataran-tataran. Ing
salajengipun kita bakal nampa prinsip iki. Mula kita uga
bakal nampa Aksioma Pasangan.

Kanthi aksioma anyar iki, kita bisa mbuktekake anane akeh
himpunan liyane. Upamane:

\begin{prop}\ollabel{prop:pairsconsequences}
Kanggo sembarang himpunan $a$ lan $b$, himpunan iki ana:
	\begin{enumerate}
		\item\ollabel{singleton} $\{a\}$
		\item\ollabel{binunion} $a \cup b$
		\item\ollabel{tuples} $\tuple{a, b}$
	\end{enumerate}
\end{prop}

\begin{proof}
\olref{singleton}. Miturut Pasangan, $\{a, a\}$ ana,
lan miturut Ekstensionalitas iku padha karo $\{a\}$.

\olref{binunion}. Miturut Pasangan, $\{a, b\}$ ana.
Saiki $a \cup b = \bigcup \{a, b\}$ ana miturut Gabungan.

\olref{tuples}. Miturut \olref{singleton}, $\{a\}$ ana.
Miturut Pasangan, $\{a, b\}$ ana. Saiki
$\{\{a\}, \{a, b\}\} = \tuple{a, b}$ ana, kanthi
ngetrapake Pasangan maneh.
\end{proof}

\begin{prob}
Tuduhna yen kanggo sembarang himpunan $a, b, c$,
himpunan $\{a, b, c\}$ ana.
\end{prob}

\begin{prob}
Tuduhna yen kanggo sembarang himpunan $a_1, \ldots, a_n$,
himpunan $\{a_1, \ldots, a_n\}$ ana.
\end{prob}

%\begin{proof} Kanthi Pasangan kaping pindho, $\{\{a_1, a_2\}, \{a_1,
%   a_3\}\}$ ana. Miturut Gabungan lan Ekstensionalitas, $\bigcup \{\{a_1,
%   a_2\}, \{a_1, a_3\}\} = \{a_1, a_2, a_3\}$ ana. Baleni cara iki
%   sakperlu. \end{proof}

\end{document}
